\documentclass[10pt,openany]{book}
\usepackage{robbook}
\usepackage[
  contentBlocks,
  fencedCode,
  pipeTables,
  relativeReferences,
  smartEllipses,
  strikeThrough,
  tableCaptions,
  taskLists,
  texMathDollars,
  tightLists=false
]{markdown}

\SetROBAccent{ROBAmber}
\SetROBVolume{VOLUME 8 - SWIFT AND OBJECTIVE-C DEVELOPERS}{ENGINEERING CEREBRO}
\SetROBChapterPrefix{Chapter}
\lstdefinelanguage{text}{}
\lstdefinelanguage{swift}{}
\lstdefinelanguage{objc}{}
\lstdefinelanguage{json}{}

\hypersetup{
  pdftitle={Building R.O.B. Volume 8: Engineering Cerebro},
  pdfauthor={Rodolfo Aramayo / Orbitus Robotics},
  pdfsubject={An implementation-level Swift and Objective-C guide to Cerebro hardware IO, three-camera perception, local and cloud models, consent-based face identity, private Messages, recording, video, and bounded autonomy},
  pdfkeywords={Swift, Objective-C, macOS, serial, AVFoundation, Vision, SceneKit, MLX, Gemini, Luxonis, AdaFace, face identity, Messages, recording, Insta360, H.264, robotics}
}

\begin{document}

\ROBTitle{Building R.O.B. - Volume 8}{ENGINEERING CEREBRO}{Three-camera perception, private AI, face identity, Messages, and recorded evidence}{An implementation-level guide for Swift and Objective-C developers}{cerebro-perception-control-lesson.png}

\ROBFrontMatter{Build a robot mind from bounded, inspectable parts}{Cerebro becomes reliable when every camera frame, serial byte, model turn, and actuator request crosses an explicit contract.}{%
This volume reads the current Cerebro repository as a living mixed-language robot system. It dissects ROBSerialBox and its POSIX/IOKit device paths; separate face, belly, and Insta360 camera roles; Vision and SceneKit rendering; MLX and Gemini inference; consent-based encrypted face identity with selectable AdaFace encoders; the fail-closed Messages bridge and optional encrypted memory; explicit synchronized recording; SceneSnapshot; bounded autonomy; safe shows; and the three-feed Vision Pro contract taught in Volume 7.

Paths are relative to the public \url{https://github.com/RudyAramayo/Cerebro} repository, with the spatial client in \url{https://github.com/RudyAramayo/ROBControllerVision}. Every chapter names the exact implementation or test being analyzed. Historical artifacts are not automatically recommended, model output is never motor authority, and physical behavior must be verified under a controlled commissioning procedure.
}

\begin{safetybox}[title={SOFTWARE CONFIDENCE IS NOT PHYSICAL CERTIFICATION}]
Cerebro controls high-energy mobile and articulated hardware. A successful build, camera detection, model response, serial banner, green UI state, or passing fixture does not certify wiring, brakes, limits, calibration, collision avoidance, force, stability, or emergency stopping. Test unpowered first, isolate actuators, use conservative limits, maintain a physical emergency stop, and keep one responsible human in command.
\end{safetybox}

\ROBSeriesMap{Volume 8 is the robot-side companion to Volume 7. Read Volume 7 for ROBControllerVision's Swift 6 state model, explicit controller ownership, dual-arm input, three H.264 feeds, immersive 360 presentation, and speech. Read this volume for Cerebro's matching server, camera, identity, private-conversation, recording, model, hardware, and safety architecture.}
\ROBSourceRoot

\begin{figure}[p]
\centering
\includegraphics[width=0.88\textwidth,height=0.76\textheight,keepaspectratio]{cerebro-perception-control-lesson.png}
\caption{Conceptual illustration: Cerebro admits camera, depth, serial, local-model, cloud-model, show, and Vision Pro data through separate bounded paths. It is a teaching diagram, not a wiring plan or literal screenshot.}
\end{figure}
\clearpage

\tableofcontents
\clearpage
\mainmatter
\setlength{\emergencystretch}{4em}
\sloppy

\markdownInput{volume-8-engineering-cerebro.md}

\end{document}
